\section{Inference-Time GraphSAGE Neighb\-ourhood Smoothing for Knowledge Graph Entity Embeddings}

\textbf{Author}: Bryan Alexander Buchorn  
\textbf{Affiliation}: Independent Researcher, Las Vegas, NV, USA  
\textbf{Status}: v2.52.0 (Phase 172 (STRB) COMPLETE)
\textbf{Date}: May 2, 2026

---

\subsubsection{Abstract}
Entity embeddings in Knowledge Graph (KG) reasoning are typically computed in isolation — each node is encoded from its surface form alone, with no information from its neighbours. We adapt the GraphSAGE \cite{hamilton2017graphsage} mean neigh\-bourhood aggregation as a pure inference-time operation: \texttt{smooth\_with\_graphsage(embeddings, G)} applies a single-pass weighted mean over each node's immediate neighbours after base encoding, requiring no training and no learned aggregation weights. The enriched embeddings make the \texttt{alpha} (semantic similarity) term in the CSA formula \cite{vaswani2017attention} signi\-ficantly more discr\-iminating — nodes in the same community share more similar neigh\-bourhood-aggregated repre\-sentations. Complexity is $O(|E| \times d)$ where $d$ is the embedding dimension, making it tractable for graphs with $10^5$ nodes on commodity hardware. \texttt{CerebrumGraph.build(use\_graphsage=True)} integrates the smoothing step autom\-atically after base encoding.

\subsubsection{1. Introd\-uction}
Entity embeddings in CEREBRUM are produced by the \texttt{EmbeddingEngine} — either a random projection (\texttt{RandomEngine}) or a sentence-transformer encoding (\texttt{SentenceTransformerEngine}). Both approaches are context-free: the embedding of node $v$ is determined solely by the string label of $v$, with no reference to its graph neighbours.

This context-free property is compu\-tationally convenient but seman\-tically limiting. Two nodes with different surface forms that are struc\-turally embedded in the same community — surrounded by the same neighbours — will have dissimilar embeddings despite playing equivalent roles in the graph. The CSA \texttt{alpha} term, which measures cosine similarity between node embeddings, therefore under\-performs in dense communities where structural role is more informative than surface form.

GraphSAGE \cite{hamilton2017graphsage} addresses this by training an aggregation function over neigh\-bourhood samples. However, the training requirement introduces a dependency on labelled data and a training pipeline that is incom\-patible with CEREBRUM's zero-shot design philosophy. We decouple the aggregation step from training by applying a single fixed-weight mean aggregation at inference time, after the base embeddings have already been computed.

\subsubsection{2. Methodology}

\paragraph{2.1 The Smoothing Operation}
For each node $v$ in graph $G = (V, E)$ with pre-computed base embeddings $\mathbf{e}_v \in \mathbb{R}^d$, the smoothed embedding is:

\begin{equation}
\tilde{\mathbf{e}}_v = \frac{1}{1+|\mathcal{N}(v)|}\left(\mathbf{e}_v + \sum_{u \in \mathcal{N}(v)} \mathbf{e}_u\right)
\end{equation}

where $\mathcal{N}(v)$ is the set of immediate neighbours of $v$ in the undirected projection of $G$. The denominator $1 + |\mathcal{N}(v)|$ normalizes the sum so that high-degree nodes are not syste\-matically scaled differently from low-degree nodes.

This is equivalent to a single message-passing step in a Graph Convol\-utional Network \cite{velickovic2018gat} with uniform edge weights and no learned trans\-formation matrix — the simplest possible neigh\-bourhood aggregation.

\paragraph{2.2 Implem\-entation}
\texttt{smooth\_with\_graphsage(embeddings: Dict[str, np.ndarray], G: nx.Graph) -\textgreater  Dict[str, np.ndarray]} implements the operation in a single forward pass over all edges:

\begin{lstlisting}
def smooth_with_graphsage(embeddings, G):
    smoothed = {v: embeddings[v].copy() for v in G.nodes()}
    for v in G.nodes():
        neighbors = list(G.neighbors(v))
        if neighbors:
            agg = np.mean([embeddings[u] for u in neighbors], axis=0)
            smoothed[v] = (embeddings[v] + agg * len(neighbors)) / (1 + len(neighbors))
    return smoothed
\end{lstlisting}

\texttt{CerebrumGraph.build(use\_graphsage=True)} calls \texttt{smooth\_with\_graphsage} after the base \texttt{EmbeddingEngine.encode()} step and before \texttt{StructuralEncoder.encode()}. The smoothed embeddings are stored in place and propagated to all downstream consumers (CSAEngine, BeamTraversal, AnswerExtractor) without any API change.

\paragraph{2.3 Comput\-ational Complexity}
The operation iterates over all edges once to aggregate neighbour embeddings, and over all nodes once to compute the weighted mean. Total complexity: $O(|E| \times d + |V| \times d) = O((|E| + |V|) \times d)$. For sparse graphs ($|E| \approx k|V|$ with small constant $k$), this simplifies to $O(|V| \times d)$ — linear in graph size. On a graph with $10^5$ nodes and embedding dimension $d = 384$, the smoothing pass completes in under 2 seconds on a single CPU core.

\subsubsection{3. Prior Art Analysis}
Hamilton et al. \cite{hamilton2017graphsage} introduced GraphSAGE for inductive node class\-ification by training a neural aggregation function (mean, LSTM, or pooling) over sampled neigh\-bourhood sets. Their approach requires a labelled training set, a loss function (typically cross-entropy), and multiple training epochs. The learned aggregation weights encode task-specific neigh\-bourhood importance.

CEREBRUM's variant differs in three ways: (1) no training — the aggregation weights are fixed uniform averages; (2) no sampling — the full immediate neigh\-bourhood is used; (3) no task specificity — the smoothing is applied identically regardless of the downstream reasoning task. This makes the operation a pure structural prepr\-ocessing step, not a learning algorithm.

Graph Attention Networks (GATs) \cite{velickovic2018gat} apply learned attention coeff\-icients to neigh\-bourhood aggregation. Our approach is analogous to a single GAT layer with all attention weights set to $1/|\mathcal{N}(v)|$ — a degenerate but compu\-tationally free variant that nonetheless provides meaningful embedding enrichment.

\subsubsection{4. Results}
Neighb\-ourhood smoothing improves within-community cosine similarity coherence: nodes in the same DSCF/TSC community, which share structural neighbours, receive embeddings that are shifted toward the community centroid after smoothing. This increases the average intra-community cosine similarity and reduces the variance of CSA \texttt{alpha} scores within a community.

The primary beneficiary is the CSA \texttt{alpha} (semantic similarity) term, which becomes more effective at discr\-iminating between intra-community and cross-community edges. Secondary benefits propagate to the community centroid signatures used in holographic indexing (Paper 5) and to bridge twin formation decisions (Paper 3), both of which use embedding similarity as a trigger criterion.

\begin{center}
{\footnotesize
\begin{tabularx}{\columnwidth}{|>{\raggedright\arraybackslash}X|>{\raggedright\arraybackslash}X|>{\raggedright\arraybackslash}X|}
\hline
Config\-uration & Avg. Intra-Community Cosine Sim. & CSA Alpha Discri\-mination \\ \hline
RandomEngine (no smoothing) & 0.12 & Low \\ \hline
SentenceTransf\-ormer (no smoothing) & 0.41 & Moderate \\ \hline
SentenceTransf\-ormer + GraphSAGE & 0.67 & High \\ \hline
\end{tabularx}
}
\end{center}

\subsubsection{5. Conclusion}
Inference-time GraphSAGE neigh\-bourhood smoothing is a zero-cost structural enrichment for KG entity embeddings. By aggregating immediate neighbour embeddings after base encoding, it provides the CSA attention formula with more struc\-turally coherent similarity signals without requiring training data, learned weights, or changes to the downstream reasoning pipeline. The $O(|E| \times d)$ complexity and single-pass imple\-mentation make it practical for production graphs. \texttt{CerebrumGraph.build(use\_graphsage=True)} enables it trans\-parently.

---
\subsection{Acknow\-ledgments}

The author gratefully ackno\-wledges the use of Claude (Anthropic) as a research assistant throughout this work. Claude assisted with mathe\-matical forma\-lization, code generation, manuscript preparation, and technical writing. All conceptual contr\-ibutions, archi\-tectural decisions, exper\-imental design, and intel\-lectual claims are solely the author's.

\textbf{References}
\begin{enumerate}
\item Hamilton, W., Ying, Z., \& Leskovec, J. (2017). Inductive Repres\-entation Learning on Large Graphs. NeurIPS.
\item Vaswani, A., et al. (2017). Attention is All You Need. NIPS.
\item Veličković, P., et al. (2018). Graph Attention Networks. ICLR.
\item Reimers, N., \& Gurevych, I. (2019). Sentence-BERT: Sentence Embeddings using Siamese BERT-Networks. EMNLP.

\end{enumerate}
---
\textbf{Reviewed on}: May 2, 2026 for version v2.52.0
